% !TEX root = ../main.tex
\begin{abstract}
A policy's rank on nuScenes' open-loop metrics or NAVSIM's leaderboard score predicts its rank at a
new deployment site only about as well as a coin flip ($\rho=-0.54$--$+0.49$ on near-pedestrian
scenes): a benchmark score reports an outcome, not the behaviour behind it.

%120
Attempting to predict a policy's performance at a new deployment site from an existing leaderboard
therefore works poorly, because a rank reflects only the outcome, not the behaviour behind it.

We therefore need a method that reveals a policy's true response, and the part of that response which
is common across situations.

We propose a \emph{model check-up}: it diagnoses a policy's true response
through pixel-level counterfactual edits, distils causal feature axes that are invariant across
domains, and thereby predicts how different models will be graded at the deployment site, lowering
the testing cost and risk of real deployment. 

By constructing a left-hand-drive to right-hand-drive
transfer experiment on scenes in which pedestrians appear, we show that the check-up's assessment on
the right-hand-drive deployment side agrees with true performance at $\rho={\NumSpBeh}$, which
validates the effectiveness of the check-up.

\end{abstract}
%120
% 此时如果我们想依据已有排名表来预测在新部署场景下的策略表现，往往会发现效果不佳，因为排名只反映了结果，而非策略背后的行为。
% 因此我们需要一种方法来揭示策略在不同情境下的具备共性的真实反应。
% 我们提出了一种“模型体检”方法，旨在通过像素级的反事实编辑来诊断策略的真实反应，提炼出跨域不变的因果特征轴，以有效预测不同模型在部署场景下的表现评级。降低实际部署中的测试成本和风险。

% 我们通过构建行人出没特定场景下的左舵测试到右舵迁移的实验，证明了模型体检方法在右舵部署侧的评估和真实表现接近率在  
% $\rho={\NumSpBeh}$，验证了模型体检方法的有效性。